% ============================================================
% PredictionRTS 项目总结 —— 导师汇报用
% 编译: XeLaTeX 或 pdfLaTeX (需 ctex 宏包)
% 图片路径: 相对本文件，output 使用 ../../output/
% 自备图请放入 figures/ 子目录并修改下方路径
% ============================================================
\documentclass[UTF8,11pt]{ctexart}
\usepackage{lmodern}
\usepackage{amsmath,amssymb}
\usepackage{graphicx}
\usepackage{booktabs}
\usepackage{array}
\usepackage{xcolor}
\usepackage{tikz}
\usepackage{float}
\usepackage{subcaption}
\usepackage{caption}
\usepackage{hyperref}
\usepackage{geometry}
\usepackage{enumitem}
\usepackage{multirow}

\geometry{top=2cm,bottom=2cm,left=2cm,right=2cm}
\definecolor{darkblue}{HTML}{0044FF}
\definecolor{darkred}{HTML}{8B0000}

% 图注与表格题
\captionsetup{font=small,labelfont=bf}
\setlength{\parindent}{2em}

\title{\textbf{PredictionRTS 项目总结}\\{\large 功能、方法创新与实验结果}}
\author{汇报人：\qquad\qquad\quad 日期：\today}
\date{}

\begin{document}
\maketitle
\vspace{-1em}
\noindent\rule{\linewidth}{0.8pt}
\vspace{1em}

% ==================== 一、背景、动机与主要问题 ====================
\section{背景、动机与主要问题}
\label{sec:motivation-problem}

\textbf{背景。} 即时战略游戏（RTS）状态高维、时序长、奖励稀疏，传统强化学习（价值函数或策略梯度）在此类设定下样本效率与长程依赖建模均面临挑战。另一方面，智能体与环境交互产生的轨迹本质上是「状态—动作—奖励」的序列，若将\textbf{轨迹预测}与\textbf{决策}视为序列建模问题，则可与 Trajectory Transformer、Decision Transformer 等序列建模范式对接，便于利用离线数据并提升可解释性。

\textbf{动机。} 在离线采集的 RTS 轨迹上，我们希望不仅学习“下一状态/下一动作”的统计规律，更希望利用状态空间已有的\textbf{几何结构}（如：哪些状态在战局语义上相近、哪些轨迹在整体走势上相似）。若能将这种结构显式注入模型——例如通过状态距离矩阵约束嵌入空间或注意力——有望提升预测一致性、抑制不合理跳变，并为策略分析与决策推荐提供可解释的几何视角。

\textbf{本项目解决的主要问题。} 在 RTS 离线轨迹场景下，\textbf{如何定义与计算状态/轨迹的几何结构，并将其系统性地转化为 Transformer 类模型的归纳偏置}，从而在状态序列预测与动作（决策）预测两项任务上同时提升效果与可解释性。具体包括：状态空间的离散化与度量、距离矩阵的构建与使用、以及“距离诱导”的嵌入初始化、注意力偏置与解码先验的设计与验证。

\vspace{0.5em}

% ==================== 二、项目主要功能 ====================
\section{项目主要功能}
\label{sec:function}

为应对上述问题，本项目实现从数据采集、状态抽象、距离建模到序列/决策模型训练与评估的完整流程，主要包含三部分功能。

\subsection{数据采集与结构化}
在 StarCraft II（PySC2）环境中运行分层 Q-Learning 智能体（SmartAgent），与内置 Bot 对战并采集：
\begin{itemize}[leftmargin=*]
    \item \textbf{状态序列}（state\_log）：每局轨迹对应的状态节点 ID 序列；
    \item \textbf{动作序列}（action\_log）：聚类策略与作战动作的离散符号序列；
    \item \textbf{奖励序列}（r\_log）：由每步净生命值变化得到的即时奖励；
    \item \textbf{对局结果}（game\_result）：胜/负/平及步数、得分等。
\end{itemize}
原始状态通过影响图哈希离散化，再经 BK-Tree 与聚类组织为可索引的状态空间，为后续距离计算与模型输入提供统一表示。

\subsection{轨迹与决策建模}
\begin{itemize}[leftmargin=*]
    \item \textbf{状态序列预测}：以 Trajectory Transformer 或 SGTransformer 对状态索引序列做自回归建模，给定前 $N$ 步预测后 $K$ 步，支持 MDS 嵌入初始化、相似度引导注意力（SG）及解码阶段的空间先验；
    \item \textbf{决策预测}：以 Decision Transformer（DT）对「回报—状态—动作」三元组建模，根据历史状态与回报预测下一步动作，用于离线策略表达与动作准确率评估。
\end{itemize}

\subsection{分析与可视化}
基于自定义状态距离与 DTW 构建状态距离矩阵与轨迹距离矩阵，进而生成适应度景观（Fitness Landscape）与状态价值地形（State Landscape）；训练与评估结果导出为 Excel 表格及 PDF/PNG 图表，便于消融与案例展示。

\vspace{0.5em}

% ==================== 三、主要难点 ====================
\section{主要难点}
\label{sec:difficulty}

\begin{enumerate}[leftmargin=*]
    \item \textbf{状态空间离散化与语义一致性}：RTS 原始观测高维、连续，需在保留战局语义的前提下映射为离散状态 ID，并保证“相近状态得到相近 ID/表征”，否则后续距离与 Transformer 建模难以生效。
    \item \textbf{状态/轨迹距离的定义与计算}：需设计兼顾单位坐标、生命值、敌我结构的度量，并在 BK-Tree 与聚类结构上高效计算两两距离及 DTW，计算量与缓存策略需权衡。
    \item \textbf{几何先验注入模型}：如何将状态距离矩阵转化为 Transformer 的归纳偏置（MDS 初始化、注意力偏置、解码先验），使模型“看见”状态空间拓扑，同时不过度约束数据驱动学习。
    \item \textbf{状态序列与决策任务的衔接}：状态预测与动作预测既可分开建模，也可在 DT 中统一；需明确评估指标（按步距离误差、动作准确率、与 Fitness 的相关性）以体现方法价值。
\end{enumerate}

\vspace{0.5em}

% ==================== 四、动机与方法创新 ====================
\section{动机与方法创新}
\label{sec:motivation}

\subsection{研究动机}
针对第一节所述问题，本项目的思路是：将轨迹与决策统一为序列建模，并\textbf{显式利用状态/轨迹的几何结构}（距离矩阵）作为模型的归纳偏置，以提升预测一致性与决策可解释性。

\subsection{方法创新点}
\begin{itemize}[leftmargin=*]
    \item \textbf{状态离散化与几何结构}：通过 BK-Tree 与聚类得到离散状态 ID，再结合自定义状态距离（坐标+生命值+匈牙利匹配）与 DTW 定义轨迹间距离，使“相近轨迹”在度量空间中有明确含义。
    \item \textbf{距离诱导的 Transformer 表征}：
    \begin{itemize}
        \item MDS 初始化：用状态距离矩阵做多维尺度变换，将得到的低维坐标作为 Embedding 或 DT 的 state embedding 初始值，实现隐空间与物理/逻辑距离的同构；
        \item SGTransformer：在自注意力中引入由距离矩阵导出的相似度偏置（sim\_bias），可学习缩放，使注意力更关注几何近邻状态；
        \item 空间先验：自回归解码时用当前状态与候选状态的距离构造高斯先验，与 logits 相乘后再 argmax，抑制几何上不合理的跳变。
    \end{itemize}
    \item \textbf{RTS 场景下的 Decision Transformer}：将 DT 应用于离线 RTS 轨迹（状态=状态 ID，动作=离散动作表，RTG 由每步奖励反向累加），并用 MDS 初始化 state embedding，在「状态+回报$\to$动作」框架下评估动作预测准确率及与 Fitness 的关系。
    \item \textbf{适应度/状态景观可视化}：基于轨迹 DTW 距离与对局结果生成 Fitness Landscape；基于状态距离矩阵与状态价值生成 State Landscape，用于定性分析预测轨迹在价值地形上的分布。
\end{itemize}

\vspace{0.5em}

% ==================== 五、核心方法 ====================
\section{核心方法}
\label{sec:methods}

\subsection{数据与距离管线}
\begin{itemize}[leftmargin=*]
    \item \textbf{BK-Tree}：按聚类组织状态节点，配合自定义距离实现快速最近邻与层次化状态索引。
    \item \textbf{状态距离矩阵}：对任意两状态，从 BK-Tree 取原始状态表示，经 \texttt{custom\_distance}（坐标+生命值等多维度量）得到标量或向量距离，再转为欧氏范数，形成方阵 $\mathbf{D}_{\mathrm{state}}$。
    \item \textbf{轨迹 DTW 距离}：以 $\mathbf{D}_{\mathrm{state}}$ 为局部代价，对轨迹对做 DTW，得到轨迹间距离矩阵，用于 Fitness Landscape 的 MDS 降维与插值。
\end{itemize}

\subsection{模型结构概览}
\begin{itemize}[leftmargin=*]
    \item \textbf{Trajectory Transformer}：Embedding + 可学习位置编码 + Transformer Encoder（因果掩码），输出头预测下一 token（状态 ID），损失为交叉熵（忽略 PAD）。
    \item \textbf{SGTransformer}：与上述结构类似，但每层自注意力在因果 mask 上叠加外部 \texttt{sim\_bias}（由 $\mathbf{D}_{\mathrm{state}}$ 等构造），并带可学习缩放因子。
    \item \textbf{Decision Transformer}：对 RTG、状态 ID、动作 ID 分别嵌入并加时间步嵌入，按 $[R_1,S_1,A_1,R_2,S_2,A_2,\ldots]$ 拼接，经 Transformer Encoder 后取状态位置的输出预测动作 logits，损失为动作交叉熵。
\end{itemize}

\subsection{训练与评估}
\begin{itemize}[leftmargin=*]
    \item 状态序列：AdamW，可选 MDS 初始化/冻结或可训练、空间先验、SG；按步预测误差在测试集上按步取平均，并在地形图上可视化输入/真实/预测轨迹。
    \item DT：AdamW，可选 MDS 初始化 state embedding；评估指标为动作预测准确率及与轨迹 Fitness 的关联（散点图、分布图）。
\end{itemize}

\vspace{0.5em}

% ==================== 六、实验结果展示 ====================
\section{实验结果展示}
\label{sec:experiments}

本节预留图表位置，请将相应图片放入 \texttt{figures/} 目录（或修改路径为 \texttt{../../output/} 下已有文件），并视需要更新图注。

\subsection{系统/方法架构图（建议放置处）}
图~\ref{fig:arch} 用于展示从原始数据到模型输出的整体流程（PySC2 $\to$ BK-Tree $\to$ 距离矩阵 $\to$ Transformer/DT $\to$ 指标）。\textbf{可替换为根据附录提示词生成的架构图。}

\begin{figure}[H]
    \centering
    % 替换为您的架构图路径，例如： figures/architecture.pdf 或 ../../output/figures/arch.png
    \fbox{\parbox{0.85\linewidth}{\centering\vspace{4cm}\textbf{[此处插入：项目整体架构图]}\\建议尺寸：宽度约 0.85\linewidth，PDF 或 PNG}}
    \caption{项目整体数据流与模型架构示意。}
    \label{fig:arch}
\end{figure}

\subsection{状态序列预测：训练曲线与消融}
图~\ref{fig:training} 展示状态序列模型的训练损失与测试集平均误差（及按步误差），对比 Vanilla / MDS 冻结（IEF）/ MDS 可训练（IET）等设置。

\begin{figure}[H]
    \centering
    % 使用项目已有图时取消下面两行注释并注释掉下面的 fbox：
    % \includegraphics[width=0.48\linewidth]{../../plot/training_comparison.png}
    % \includegraphics[width=0.48\linewidth]{../../plot/prediction_results.png}
    \begin{minipage}{0.48\linewidth}
        \centering
        \fbox{\parbox{\linewidth}{\centering\vspace{2.8cm}\textbf{[左：训练 Loss / 总平均误差]}}}
    \end{minipage}
    \hfill
    \begin{minipage}{0.48\linewidth}
        \centering
        \fbox{\parbox{\linewidth}{\centering\vspace{2.8cm}\textbf{[右：消融/预测误差对比]}}}
    \end{minipage}
    \caption{状态序列预测：训练曲线与消融实验（可替换为 \texttt{plot/training\_comparison.png} 与 \texttt{plot/prediction\_results.png} 或对应 PDF）。}
    \label{fig:training}
\end{figure}

\subsection{状态序列预测：地形上的轨迹对比}
图~\ref{fig:traj} 展示若干测试样本在状态价值地形上的输入序列、真实后续轨迹与模型预测轨迹（不同线型/颜色）。\textbf{可替换为 \texttt{output/figures/pred\_FL*/traj\_analysis\_*.pdf} 的三列对比。}

\begin{figure}[H]
    \centering
    \begin{minipage}{0.32\linewidth}
        \centering
        \fbox{\parbox{\linewidth}{\centering\vspace{2.5cm} Vanilla}}
    \end{minipage}
    \hfill
    \begin{minipage}{0.32\linewidth}
        \centering
        \fbox{\parbox{\linewidth}{\centering\vspace{2.5cm} W/ IEF}}
    \end{minipage}
    \hfill
    \begin{minipage}{0.32\linewidth}
        \centering
        \fbox{\parbox{\linewidth}{\centering\vspace{2.5cm} W/ IET}}
    \end{minipage}
    \caption{某测试样本在状态地形上的预测对比（左：Vanilla；中：MDS 冻结；右：MDS 可训练）。}
    \label{fig:traj}
\end{figure}

\subsection{Decision Transformer 评估}
图~\ref{fig:dt} 展示 DT 的动作预测准确率及其与轨迹 Fitness 的关系（如 Fitness vs. Accuracy 散点、准确率分布）。\textbf{可替换为 \texttt{output/dt\_performance\_analysis.png}、\texttt{output/dt\_results.png}。}

\begin{figure}[H]
    \centering
    % 使用项目已有图时取消注释（二选一或并排）：
    % \includegraphics[width=0.9\linewidth]{../../output/dt_performance_analysis.png}
    % \includegraphics[width=0.9\linewidth]{../../output/dt_results.png}
    \fbox{\parbox{0.9\linewidth}{\centering\vspace{3.2cm}\textbf{[此处插入：DT 性能分析图]}\\例如 Fitness vs Accuracy、准确率分布等}}
    \caption{Decision Transformer 动作预测准确率与 Fitness 关联分析。}
    \label{fig:dt}
\end{figure}

\subsection{Embedding 可视化（可选）}
若需展示 MDS 初始化或训练前后嵌入在二维上的分布，可在图~\ref{fig:embed} 位置插入 \texttt{output/embedding/*/train\_epoch\_*.pdf} 中的代表性子图。

\begin{figure}[H]
    \centering
    \fbox{\parbox{0.75\linewidth}{\centering\vspace{2.8cm}\textbf{[可选：Embedding 二维可视化]}}}
    \caption{状态嵌入在训练过程中的二维投影（可选）。}
    \label{fig:embed}
\end{figure}

\vspace{0.5em}

% ==================== 附录：文生图提示词 ====================
\appendix
\section{科研绘图用文生图提示词（Nano Banana Pro）}
\label{sec:prompts}

以下提示词适用于生成\textbf{项目架构图}与\textbf{方法机理图}，风格为科研/技术示意图：简洁、线框与框图为主、少量配色、适合论文或汇报幻灯片。

\subsection{项目整体架构图}
\begin{quote}
\textbf{英文提示词（推荐直接使用）：}

Scientific diagram, research pipeline, flowchart style. From left to right: (1) ``Raw Game'' box with small icons of units and map; (2) ``State Discretization'' with a tree icon and ``BK-Tree''; (3) ``Distance Matrix'' as a symmetric grid/heatmap; (4) ``Transformer / DT'' as a stack of layers with ``Attention'' label; (5) ``Metrics \& Visualization'' with a small plot icon. Arrows connect each block. Clean white background, thin black or dark gray lines, two accent colors (e.g. blue and orange) for key boxes. No 3D, flat 2D, suitable for academic paper. Label each block in sans-serif font.
\end{quote}

\begin{quote}
\textbf{中文补充说明（便于微调）：}

科研风格流程图，从左到右依次：原始对局数据 $\to$ 状态离散化（BK-Tree/聚类）$\to$ 状态距离矩阵（对称矩阵示意）$\to$ Transformer/Decision Transformer 模型 $\to$ 指标与可视化。箭头连接，白底、细线、少量蓝/橙强调，平面二维，适合论文插图。
\end{quote}

\subsection{方法机理图：距离诱导表征}
\begin{quote}
\textbf{英文提示词：}

Concept diagram, embedding space. Left: ``State Distance Matrix'' shown as a 2D grid or heatmap with a few labeled nodes. Center: ``MDS'' box with arrow ``low-dim coordinates''. Right: ``Transformer Embedding Layer'' as a row of small squares/cells with ``init from MDS''. One or two curved arrows from matrix to embedding to show ``geometry preserved''. Minimal style, white background, blue/gray tones, academic illustration.
\end{quote}

\begin{quote}
\textbf{中文补充说明：}

概念图：左侧状态距离矩阵（网格或热力图），中间 MDS 框与“低维坐标”箭头，右侧 Transformer 嵌入层（小方格），并标注“由 MDS 初始化”“几何保持”。简洁、白底、蓝灰色调。
\end{quote}

\subsection{方法机理图：相似度引导注意力（SG）}
\begin{quote}
\textbf{英文提示词：}

Schematic of attention mechanism. A ``Query'' and ``Key'' block, then a ``Similarity / Distance Bias'' block (e.g. from state distance matrix). Arrows: bias is added to attention scores before softmax. Below: ``Causal mask + Bias''. Flat diagram, boxes and arrows only, one accent color for the bias path, rest black/gray. Paper-style figure.
\end{quote}

\begin{quote}
\textbf{中文补充说明：}

注意力机制示意图：Query/Key、相似度/距离偏置（来自状态距离矩阵）、偏置在 softmax 前加到注意力分数上；标注“因果掩码 + 偏置”。框图+箭头，偏置路径用一强调色，其余黑灰。
\end{quote}

\subsection{方法机理图：空间先验解码}
\begin{quote}
\textbf{英文提示词：}

Decoding step diagram. Top: ``Current state'' node. Middle: ``Logits from Transformer'' and ``Spatial prior (Gaussian from distance)''; arrow ``multiply'' then ``argmax''. Bottom: ``Next state''. Simple flowchart, white background, minimal labels, academic style.
\end{quote}

\begin{quote}
\textbf{中文补充说明：}

解码步骤：当前状态 $\to$ Transformer logits 与空间先验（由距离构造的高斯）$\to$ 相乘后 argmax $\to$ 下一状态。简单流程图，白底、学术风格。
\end{quote}

\vspace{1em}
\noindent 使用建议：将上述英文提示词复制到 Nano Banana Pro 中生成图像，若需更“像论文图”，可在提示词末尾加 \texttt{, line art, no shading, vector style} 或 \texttt{, conference paper figure style} 等短语微调。

\end{document}
